% strong/topology_qh3.tex
% Part of the Hopfion Framework Repository
% Source: Paper XVI — Quark Generation Masses in the Density-Feedback Hopfion:
%   A Survey of Ruled-Out Mechanisms and the Q_H=3 Gradient-Flow Landscape
% Status: published
% Last updated: 2026-06-28
%
%% hopfion-framework/preamble.tex
%% Shared preamble for all repository files.
%% \input this at the top of any standalone document.
%% Repository files are designed to be \input into papers directly.

\usepackage{amsmath,amssymb,amsthm}
\usepackage{hyperref}
\usepackage{enumitem}
\usepackage{booktabs}
\usepackage{xcolor}
\usepackage{geometry}
\geometry{margin=2.5cm}

%── Theorem environments ──────────────────────────────────────────────────────
\newtheorem{theorem}{Theorem}[section]
\newtheorem{lemma}[theorem]{Lemma}
\newtheorem{proposition}[theorem]{Proposition}
\newtheorem{corollary}[theorem]{Corollary}
\theoremstyle{definition}
\newtheorem{definition}[theorem]{Definition}
\newtheorem{result}[theorem]{Result}
\theoremstyle{remark}
\newtheorem{remark}[theorem]{Remark}
\newtheorem{observation}[theorem]{Observation}
\newtheorem{conjecture}{Conjecture}
\newtheorem{construction}[theorem]{Construction}
\newtheorem{condition}[theorem]{Condition}

%── Repository metadata commands ─────────────────────────────────────────────
% Usage: \status{published} or \status{open} or \status{conjectured}
\newcommand{\status}[1]{\par\noindent{\color{blue}\textbf{Status:} #1}}
% Usage: \source{Paper I, Theorem thm:scale}
\newcommand{\source}[1]{\par\noindent{\color{darkgray}\textbf{Source:} #1}}
% Usage: \residual{0.013\%} or \residual{exact}
\newcommand{\residual}[1]{\par\noindent{\color{teal}\textbf{Residual:} #1}}
% Usage: \depends{foundations/pentagon, foundations/suppression}
\newcommand{\depends}[1]{\par\noindent{\color{purple}\textbf{Depends on:} #1}}

%── Framework macros ─────────────────────────────────────────────────────────
\newcommand{\vp}{\varphi}
\newcommand{\bst}{\beta^{*}}
\newcommand{\Seff}{S_{\mathrm{eff}}}
\newcommand{\TCMB}{T_{\mathrm{CMB}}}
\newcommand{\Lcond}{\Lambda_{\mathrm{cond}}}
\newcommand{\LUV}{\Lambda_{\mathrm{UV}}}
\newcommand{\LQCD}{\Lambda_{\mathrm{QCD}}}
\newcommand{\MPl}{M_{\mathrm{Pl}}}
\newcommand{\ERy}{E_{\mathrm{Ry}}}
\newcommand{\aem}{\alpha_{\mathrm{em}}}
\newcommand{\sinW}{\sin^{2}\!\theta_{W}}
\newcommand{\vEW}{v_{\mathrm{EW}}}
\newcommand{\twoI}{2I}
\newcommand{\twoT}{2T}
\newcommand{\twoO}{2O}
\newcommand{\Ehat}{\hat{E}_8}
\newcommand{\QH}{Q_H}
\newcommand{\Ztwo}{\mathbb{Z}_2}
\newcommand{\Zthree}{\mathbb{Z}_3}
\newcommand{\SU}{\mathrm{SU}}
\newcommand{\rCMB}{\rho_{\mathrm{CMB}}}
\newcommand{\Lobs}{\Lambda_{\mathrm{obs}}}
\newcommand{\rinfty}{\rho_{\infty}}
\newcommand{\oBD}{\omega_{\mathrm{BD}}}
\newcommand{\xNMC}{\xi_{\mathrm{NMC}}}
\newcommand{\Efb}{E_{\mathrm{fb}}}
\newcommand{\Kfb}{K_{\mathrm{fb}}}
\newcommand{\dimq}[1]{d_{#1}}
\newcommand{\deff}{d_{\mathrm{eff}}}
\newcommand{\Kth}{K_{\mathrm{th}}}
\newcommand{\Ja}{\mathcal{J}_a}
\newcommand{\Jfour}{\mathcal{J}_4}
\newcommand{\Jiso}{\mathcal{J}_{\mathrm{iso}}}
\newcommand{\rhoinf}{\rho_{\infty}}
\newcommand{\rhoa}{\rho_{\mathrm{atom}}}
\newcommand{\Heff}{H_{\mathrm{eff}}}
\newcommand{\epsr}{\varepsilon(r)}
\newcommand{\Hcond}{H_{\mathrm{cond}}}
\newcommand{\eps}{\varepsilon}
\newcommand{\Rzero}{R_0}
\newcommand{\Cstar}{C^*}
\newcommand{\mH}{m_H}
\newcommand{\aem}{\alpha_{\mathrm{em}}}
\newcommand{\ms}{\mu^*}
\newcommand{\lam}{\lambda}
\newcommand{\fzero}{f_0}
\newcommand{\ftwo}{f_2}
\newcommand{\Itor}[1]{I_{#1}^{\mathrm{tor}}}
\newcommand{\Jfb}{\mathcal{J}_{\mathrm{fb}}}
\newcommand{\bz}{\beta_0}
\newcommand{\sopt}{s_{\mathrm{opt}}}
\newcommand{\leff}{\lambda_{\mathrm{eff}}}
\newcommand{\ord}{\operatorname{ord}}
\newcommand{\G}{\mathcal{G}}
\newcommand{\Ia}{I_{2a}}
\newcommand{\Ib}{I_{4}}
\newcommand{\kk}{\kappa}
\newcommand{\aInv}{\alpha^{-1}}
\newcommand{\mWZW}{m_{\mathrm{WZW}}}
\newcommand{\tang}{\dot\Gamma}
\newcommand{\Jones}{V}
\newcommand{\lk}{\mathrm{lk}}
\newcommand{\phistar}{\phi^*}
\newcommand{\CC}{\mathbb{C}}
\newcommand{\HH}{\mathbb{H}}
\newcommand{\OO}{\mathbb{O}}
\newcommand{\ZZ}{\mathbb{Z}}
\newcommand{\pcond}{p_{\mathrm{cond}}}
\newcommand{\pQM}{p_{\mathrm{QM}}}
\newcommand{\xcond}{x_{\mathrm{cond}}}
\newcommand{\CHSH}{\mathrm{CHSH}}
\newcommand{\Tsirelson}{2\sqrt{2}}
\newcommand{\RR}{\mathbb{R}}
\newcommand{\Hil}{\mathcal{H}}
\newcommand{\CQ}{\mathcal{C}_Q}
\newcommand{\CQr}{\mathcal{C}_Q^{\mathrm{red}}}
\newcommand{\bOmega}{\boldsymbol{\Omega}}
\newcommand{\dmu}{d\mu}
\newcommand{\nlvl}[1]{n\!\left(#1\right)}
\newcommand{\IE}{\mathrm{IE}}
\newcommand{\bsF}{{}_2F_1}




\section*{$\QH=3$ Sector Topology and Numerical Construction}
\label{sec:topology_qh3}

% Results extracted from Paper XVI
% Scan date: 2026-06-28

%════════════════════════════════════════════════════════════════════
% String Tension and Junction Energy
%════════════════════════════════════════════════════════════════════

\begin{proposition}[The Y-junction arms lie off the trefoil curve]
\label{P16:prop:arms_off_curve}
The three Y-junction arms, each connecting the origin to a crossing
midpoint $\mathbf{m}_k$, are not subsets of the trefoil centre curve
$\Gamma(s)$. Consequently, the string tension $\sigma$ cannot be read
off from the tube integrals $J_4$ or $J_{2a}$; a genuine computation
of $\sigma$ requires the condensate field's behaviour in the bulk region
between the tube and the Fermat--Steiner point.
\end{proposition}
\status{published}
\source{P16:prop:arms_off_curve}

\begin{proposition}[Junction energy is negligible]
\label{P16:prop:ejunc_zero}
The Y-junction energy $E_\cup(C^*)$, computed by gradient-flow
minimisation of $E_{\mathrm{geom}}=\Kfb\times J_4$ in the ball
$|\mathbf{x}|<2$ centred on the Fermat--Steiner origin with Dirichlet
boundary conditions from Construction~\ref{P16:constr:perstrand}, satisfies
\begin{equation}
  \frac{E_\cup}{E_{\mathrm{baryon}}}
  \;=\; 6.9\times10^{-8}
  \quad(C^*=2.5062),
\end{equation}
negligible at $C^*=2.0$ ($1.8\times10^{-6}$) and $C^*=3.0$
($1.8\times10^{-9}$), scaling approximately as $E_\cup\propto C^{*-16}$.
\end{proposition}
\status{published}
\source{P16:prop:ejunc_zero}
\residual{exact (numerical)}

\begin{remark}[Exact baryon energy]
\label{P16:rem:ejunc_consequence}
Proposition~\ref{P16:prop:ejunc_zero} simplifies the baryon energy to
\begin{equation}
  \boxed{E_{\mathrm{baryon}} \;=\; 1.04405\cdots\times E_{\mathrm{profile}}^*,}
\end{equation}
with $E_\cup$ discarded without approximation. The sole remaining open
quantity is the absolute value of $E_{\mathrm{profile}}^*$ in physical
units, requiring the $\QH=3$ saddle finder.
\end{remark}
\status{published}
\source{P16:rem:ejunc_consequence}

\begin{proposition}[Straight-tube tension is computable without the saddle]
\label{P16:prop:sigma_tot}
The total energy per unit length of an infinite, straight Faddeev--Niemi
flux tube,
\begin{equation}
  \sigma_{\mathrm{tot}}
  \;=\; 2\pi\int_0^\infty \Kfb^{(2D)}(r)\;\rho_{J_4}^{(2D)}(r)\;r\,\mathrm{d}r,
\end{equation}
converges and evaluates to $\sigma_{\mathrm{tot}} = 5392.1$ (in natural
units with $C_3^*=2.5062$, relative error $<10^{-6}$). Under the
rescaling $u=r\,C_3^*$, the exterior-to-interior ratio
\begin{equation}
  \frac{\sigma_{\mathrm{ext}}}{\sigma_{\mathrm{int}}} = 0.04405447\ldots
\end{equation}
is a universal constant independent of $C_3^*$.
\end{proposition}
\status{published}
\source{P16:prop:sigma_tot}
\residual{exact (numerical, $<10^{-6}$)}

\begin{proposition}[The exterior string tension is analytically determined]
\label{P16:prop:sigma_ext_formula}
Paper~XV's string tension satisfies
\begin{equation}
  \sigma_{\mathrm{ext}} \cdot 3R_0
  \;=\; 0.04405447\cdots \times E_{\mathrm{profile}}^*,
\end{equation}
following purely from the profile ansatz, requiring neither the 3D saddle
nor a bulk field equation. Consequently: (i)~$\sigma_{\mathrm{ext}}\cdot3R_0$
is a generation-universal correction; (ii)~no bulk field computation is
required for $\sigma$; (iii)~the generation-dependent part of the quark
mass correction does not reside in $\sigma\cdot3R_0$.
\end{proposition}
\status{published}
\source{P16:prop:sigma_ext_formula}
\residual{exact}

\begin{remark}[Three-flow decomposition of the baryon energy]
\label{P16:rem:three_flows}
All three terms of $E_{\mathrm{baryon}} = \Kfb\times J_4 + \sigma\cdot3R_0
+ E_\cup$ are now determined: (i)~$\Kfb\times J_4$ is computable
analytically ($E_{\mathrm{profile}}^*=18/\vp^{17}$) and numerically
($K_{\min}\approx2294$ at $N=192$); (ii)~$\sigma\cdot3R_0 =
0.04405\times E_{\mathrm{profile}}^*$ analytically
(Proposition~\ref{P16:prop:sigma_ext_formula}); (iii)~$E_\cup\approx
7\times10^{-8}E_{\mathrm{baryon}}$ numerically (Proposition~\ref{P16:prop:ejunc_zero}).
\end{remark}
\status{published}
\source{P16:rem:three_flows}

\begin{remark}[Relation between $\sigma_{\mathrm{tot}}$ and Paper~XV's $\sigma$]
\label{P16:rem:sigma_decomposition}
$\sigma_{\mathrm{tot}}$ is not the same as Paper~XV's $\sigma$.
In the three-term decomposition, $\Kfb\times J_4$ counts the internal
tube energy ($\sigma_{\mathrm{int}}$ per unit length), and $\sigma\cdot3R_0$
counts the exterior tail energy ($\sigma_{\mathrm{ext}}$ per unit length).
Proposition~\ref{P16:prop:sigma_ext_formula} shows $\sigma_{\mathrm{ext}}/
\sigma_{\mathrm{int}}=0.04405$: a $4.4\%$ correction, non-negligible at
the logarithmic precision needed for quark masses.
\end{remark}
\status{published}
\source{P16:rem:sigma_decomposition}

\begin{remark}[Chern--Simons spoke for the confined baryon]
\label{P16:rem:cs_spoke_quark}
Each of the three trefoil strands carries $\QH=1$ topology individually
(Construction~\ref{P16:constr:perstrand}); the full baryon carries $\QH=3$.
The Chern--Simons spoke for the baryon is
\begin{equation}
  \Delta Q^{(q)} = Q_H^{(\text{baryon})} - Q_H^{(\text{single strand})}
  = 3 - 1 = 2,
\end{equation}
measuring the barrier that confines the three strands: a single strand
must cross $\mathcal{C}^{(q)}=4\pi\times2=8\pi$ to escape from $\QH=3$
to the free $\QH=1$ state. The leading constituent quark mass is
$m_q^{(0)} \approx 215.5$\,MeV; three quarks give
$3m_q^{(0)}\approx646.5$\,MeV, accounting for $68.9\%$ of the proton
mass $m_p=938.272$\,MeV with no free parameters.
\end{remark}
\status{published}
\source{P16:rem:cs_spoke_quark}

\begin{remark}[Junction flow: numerical results and generation-dependence conclusion]
\label{P16:rem:junction_results}
The Y-junction is a topological deconcentration locus: the three $\QH=1$
flux tubes emanating from the Fermat--Steiner point carry equal topological
charge, and their $\rho_{J_4}$ contributions cancel at the merge point by
destructive interference. The measured junction $\rho_{J_4}$ is a factor
of $\sim4\times10^6$ smaller than three straight-tube arms of the same
length would give. The generation-dependent correction to the quark mass
formula does not arise from the geometric Y-junction energy.
\end{remark}
\status{published}
\source{P16:rem:junction_results}

%════════════════════════════════════════════════════════════════════
% Phase Discontinuity in the Original Ansatz
%════════════════════════════════════════════════════════════════════

\begin{proposition}[The single-nearest-point ansatz is discontinuous at every crossing]
\label{P16:prop:phase_jump}
Let $n(\mathbf{x})$ be built from the nearest point on the trefoil
centreline with phase $\theta(\mathbf{x})=3t_{\mathrm{near}}(\mathbf{x})$.
At each of the three crossing regions, the internal phase $\theta$ jumps
by exactly $\pi$ across the surface where nearest-point ownership switches
between the over- and under-strand. This jump is not removable by mesh
refinement; it is forced by the curve's own geometry, independent of
$C^*$, $R_0$, or $r_0$.
\end{proposition}
\status{published}
\source{P16:prop:phase_jump}

\begin{remark}[The jump is not a removable gauge artifact]
\label{P16:rem:not_a_pole}
At the crossing midline, $\rho$ equals the half-separation $r_0=0.874$
between strand centrelines, giving $\sin f_0(r_0 C_3^*)=0.275$: not
small. The resulting jump $|\Delta n|=2\sin f_0(r_0 C_3^*)\approx0.55$
(out of a maximum of $2$) is a physically substantial discontinuity in
the field, not a coordinate artifact.
\end{remark}
\status{published}
\source{P16:rem:not_a_pole}

\begin{proposition}[The discontinuity carries divergent gradient energy under refinement]
\label{P16:prop:divergent_energy}
The local Faddeev--Niemi gradient energy density $g^2=\sum_a|\partial_a n|^2$
evaluated near a crossing midline on a sequence of grids of spacing $h$
has a peak value consistent with $g^2_{\max}\propto h^{-2}$ (the scaling
of a true jump discontinuity under centred finite differences), rather
than saturating at a finite value. Numerical evidence: $g^2_{\max}$
grows from $64.6$ at $h=0.0656$ to $3462.9$ at $h=0.0082$, fitting
$g^2_{\max}\propto h^{-1.92}$.
\end{proposition}
\status{published}
\source{P16:prop:divergent_energy}

\begin{remark}[A naive smoothing does not resolve the discontinuity]
\label{P16:rem:naive_smoothing_fails}
Blending the two strands' phases via a distance-weighted circular
average reduces the local gradient energy by $\approx3.1\times$ but
cannot eliminate the discontinuity: at the exact crossing midpoint the
two phases differ by $\pi$ (Proposition~\ref{P16:prop:phase_jump}), so equal
weights give the zero vector and $\arg(0)$ is undefined. This is a
generic topological obstruction: no smooth, single-phase repair can work.
Only a construction that maintains independent phase freedom for both
strands --- working in the full $S^3$ Hopf-fibre bundle rather than a
single real $\theta(\mathbf{x})$ --- can avoid the singularity.
\end{remark}
\status{published}
\source{P16:rem:naive_smoothing_fails}

\begin{remark}[Two independent mechanisms causing saddle non-convergence]
\label{P16:rem:two_mechanisms}
The phase discontinuity (Proposition~\ref{P16:prop:phase_jump}) is logically
independent of the inter-tube field-overlap mechanism identified by
Paper~XV: the overlap argument concerns the radial profile $f_0(\rho)$
and does not reference the internal phase $\theta$ at all. Both effects
may contribute additively to the observed non-convergence; this paper
establishes the phase discontinuity as a second, real, independently
verifiable candidate cause.
\end{remark}
\status{published}
\source{P16:rem:two_mechanisms}

%════════════════════════════════════════════════════════════════════
% $S^3$-Lift Construction (Local)
%════════════════════════════════════════════════════════════════════

\begin{construction}[Complex-doublet superposition ansatz]
\label{P16:constr:s3lift}
Near a crossing, build each strand's own Hopf doublet
$Z_{\mathrm{over}},Z_{\mathrm{under}}$ independently, using each strand's
own distance and fixed phase $3t_i$, then combine by
inverse-square-distance-weighted complex superposition in $\mathbb{C}^2$
before projecting to $n\in S^2$:
\begin{equation}
  Z(\mathbf{x}) = \frac{w_{\mathrm{over}}Z_{\mathrm{over}}
    + w_{\mathrm{under}}Z_{\mathrm{under}}}
    {\bigl|w_{\mathrm{over}}Z_{\mathrm{over}}
    + w_{\mathrm{under}}Z_{\mathrm{under}}\bigr|},
  \qquad w_s = d_s(\mathbf{x})^{-2}.
\end{equation}
\end{construction}
\status{published}
\source{P16:constr:s3lift}

\begin{proposition}[The complex-doublet superposition is well-defined and smooth at the crossing midline]
\label{P16:prop:s3_smooth}
At the geometric midpoint of a crossing region (equal weight,
$f_{\mathrm{over}}=f_{\mathrm{under}}\equiv f_*$), the unnormalised sum is
$Z_{\mathrm{over}}+Z_{\mathrm{under}}=(2\cos(f_*/2),\,0)$, a nonzero
vector for any $f_*\neq\pi$, normalising to $n=(0,0,1)$. The resulting
field is differentiable across the midline, in contrast to the undefined
$\arg(0)$ of the real-phase blend.
\end{proposition}
\status{published}
\source{P16:prop:s3_smooth}

\begin{theorem}[The $S^3$-lift ansatz removes the divergent gradient energy]
\label{P16:thm:s3_bounded_energy}
The local Faddeev--Niemi gradient energy density of
Construction~\ref{P16:constr:s3lift}, evaluated near crossing region~1 on a
resolution sequence from $N=41$ ($h=0.0656$) to $N=321$ ($h=0.0082$),
converges to a finite peak value near $53.1$--$53.2$ as $h\to0$, in
contrast to the single-nearest-point ansatz's divergence
($g^2_{\max}\propto h^{-1.92}$). The total local gradient energy is
similarly flat ($35.2\to45.1$, converged by $N=121$).
\end{theorem}
\status{published}
\source{P16:thm:s3_bounded_energy}

\begin{remark}[Why the complex lift succeeds where the real blend fails]
\label{P16:rem:why_s3_works}
The Hopf doublet carries an additional real degree of freedom: the
$z_1=\cos(f/2)$ component does not depend on $\theta$ and survives the
cancellation of the antipodal $e^{i\theta}$ terms. It is precisely this
$z_1$ component, invariant under the antipodal phase flip, that fixes
the combined doublet to a well-defined, nonzero value. This is a
structural feature of the full $S^3$ Hopf bundle invisible once the
doublet has been projected to a single real $S^2$ phase.
\end{remark}
\status{published}
\source{P16:rem:why_s3_works}

%════════════════════════════════════════════════════════════════════
% Global $S^3$-Lift and Hopf Charge
%════════════════════════════════════════════════════════════════════

\begin{construction}[Global two-strand $S^3$-lift ansatz]
\label{P16:constr:global_s3lift}
At each point $\mathbf{x}$, find the nearest strand parameter $t_1$
and the nearest parameter $t_2$ at curve-parameter distance $>\pi/2$
from $t_1$ (a genuinely different strand passage). Build $Z_1,Z_2$ as
in Construction~\ref{P16:constr:s3lift} using $d_1,d_2$ and phases
$\theta_i=3t_i$, and combine via inverse-square weighting. Where no
second strand is found, the construction reduces to the single-doublet
ansatz with $w_2=0$.
\end{construction}
\status{published}
\source{P16:constr:global_s3lift}

\begin{proposition}[The global ansatz reduces correctly away from crossings]
\label{P16:prop:reduces_away}
At a generic point on the trefoil tube, away from all three crossing
regions, Construction~\ref{P16:constr:global_s3lift} agrees with the original
single-nearest-point ansatz to machine precision ($6\times10^{-14}$).
\end{proposition}
\status{published}
\source{P16:prop:reduces_away}

\begin{proposition}[Global non-degeneracy]
\label{P16:prop:global_nondegen}
The unnormalised doublet $Z_1+Z_2$ (weighted as in
Construction~\ref{P16:constr:global_s3lift}) does not vanish at any sampled
point in a neighbourhood of the trefoil: a scan of $\sim6.4\times10^5$
grid points finds $|Z_1+Z_2|_{\min}=0.107$; focused scans at each
crossing find $|Z_1+Z_2|_{\min}\in\{0.386,0.386,0.465\}$.
\end{proposition}
\status{published}
\source{P16:prop:global_nondegen}

\begin{remark}[Non-degeneracy scope]
\label{P16:rem:nondegen_scope}
Propositions~\ref{P16:prop:reduces_away}--\ref{P16:prop:global_nondegen} together
show Construction~\ref{P16:constr:global_s3lift} is a well-defined, continuous
map from a neighbourhood of the trefoil into $S^2$, reducing to the
proven $\QH=3$ ansatz away from the crossings. This is a necessary
prerequisite for computing the Hopf charge, but is not itself a
computation: the charge must be computed directly.
\end{remark}
\status{published}
\source{P16:rem:nondegen_scope}

\begin{proposition}[The global $S^3$-lift ansatz is Hopf-trivial]
\label{P16:prop:s3lift_charge_zero}
The Whitehead integral evaluated on grids from $40^3$ to $80^3$
converges to $\QH\approx0$ for Construction~\ref{P16:constr:global_s3lift},
not $\QH=3$.
\end{proposition}
\status{published}
\source{P16:prop:s3lift_charge_zero}

\begin{remark}[Diagnosis: missing poloidal winding]
\label{P16:rem:missing_poloidal_winding}
The vanishing follows from a structural gap: the phase $\theta=3t$
implemented throughout has no poloidal dependence on $\chi$. The
preimage of a generic regular value consists of three small, mutually
unlinked transverse circles (one at each of three isolated curve
parameters) rather than a single connected curve with nonzero linking
number. The deficiency is not specific to the $S^3$-lift repair; it
is inherited from the underlying single-nearest-point ansatz.
\end{remark}
\status{published}
\source{P16:rem:missing_poloidal_winding}

\begin{remark}[Consequence for the saddle-finder programme]
\label{P16:rem:s3lift_not_sufficient}
Theorem~\ref{P16:thm:s3_bounded_energy}'s smoothness result stands, but
Construction~\ref{P16:constr:global_s3lift} as stated is not a valid
replacement initial condition for the saddle-finder, since it does not
carry the required $\QH=3$ topology. A corrected construction must build
in the poloidal dependence explicitly, e.g.\ $\Phi=\chi+3t$, and verify
$\QH=3$ via the same Whitehead-integral test.
\end{remark}
\status{published}
\source{P16:rem:s3lift_not_sufficient}

%════════════════════════════════════════════════════════════════════
% Profile Parameter Threshold
%════════════════════════════════════════════════════════════════════

\begin{proposition}[Critical profile parameter from the trefoil's injectivity radius]
\label{P16:prop:cstar_crit}
The trefoil curve $T_{2,3}$ has injectivity radius exactly $r_0=0.874$
(the crossings are the curve's own points of closest self-approach,
$2r_0=1.748$). This defines a critical profile parameter
\begin{equation}
  C^*_{\mathrm{crit}} = \frac{1}{r_0} \approx 1.1442,
\end{equation}
below which the field's support extends past the injectivity radius.
The physical value $C_3^*=2.5062$ gives tube radius $1/C_3^*=0.399$,
a factor of $2.2$ below $r_0$, safely in the single-nearest-point regime.
\end{proposition}
\status{published}
\source{P16:prop:cstar_crit}

\begin{remark}[Energetic signature of crossing $C^*_{\mathrm{crit}}$]
\label{P16:rem:cstar_energy_signature}
At fixed grid resolution, the total gradient energy of the
single-nearest-point ansatz shows its steepest rate of increase in the
neighbourhood $C^*\in[0.7,1.0]$, immediately below $C^*_{\mathrm{crit}}$,
before saturating for $C^*\lesssim0.3$. Three locally distinct, sharply
concentrated crossing regions are visible for $C^*>C^*_{\mathrm{crit}}$;
these merge into a single diffuse region for $C^*\llC^*_{\mathrm{crit}}$.
\end{remark}
\status{published}
\source{P16:rem:cstar_energy_signature}

\begin{proposition}[The $S^3$-lift's nonzero charge at small $C^*$ is a strand-counting artifact]
\label{P16:prop:cstar_artifact}
The Whitehead-integral charge of Construction~\ref{P16:constr:global_s3lift}
across $C^*\in[0.1,2.5062]$ is non-monotonic, peaking near
$\QH\approx-0.13$ to $-0.15$ around $C^*\approx0.3$--$0.4$ and
decaying back to zero as $C^*\to0$. Replacing the two-strand
superposition with a three-strand version reduces the peak magnitude
by roughly half ($\approx-0.05$), attributing the signal to the
two-strand truncation, not a topological effect.
\end{proposition}
\status{published}
\source{P16:prop:cstar_artifact}

\begin{remark}[Blending cannot supply the missing winding at any strand count]
\label{P16:rem:blending_insufficient}
No amount of correctly weighted strand-blending can repair the topological
deficiency: every strand contributes a phase $\theta_i=3t_i$ that is
individually poloidal-winding-free, and a sum of $\chi$-independent
functions is $\chi$-independent. A corrected construction must supply
$\chi$-dependence within each strand's own phase, e.g.\
$\Phi_i=\chi_i+3t_i$, before any crossing-region superposition.
\end{remark}
\status{published}
\source{P16:rem:blending_insufficient}

%════════════════════════════════════════════════════════════════════
% Per-Strand Construction with Poloidal Winding
%════════════════════════════════════════════════════════════════════

\begin{proposition}[The trefoil's Bishop frame has nontrivial holonomy]
\label{P16:prop:bishop_holonomy}
The Bishop (relatively parallel adapted) frame does not close up after
one traversal of $T_{2,3}$: the holonomy angle is
\begin{equation}
  H = -63.0191^\circ = -1.09989~\mathrm{rad},
\end{equation}
with exactly $\mathbb{Z}_3$-symmetric per-segment holonomy $H/3=-21.0064^\circ$.
\end{proposition}
\status{published}
\source{P16:prop:bishop_holonomy}
\residual{$<0.001^\circ$ (numerical, $N_T=80000$ samples)}

\begin{remark}[Distinction from the $\theta=3t$ winding]
\label{P16:rem:holonomy_distinct}
The holonomy $H\approx-63^\circ$ is a property of the curve's normal
bundle, unrelated to the coordinate winding number $3$ in $\theta=3t$.
The two numbers are not expected to agree: $H$ measures genuine
parallel-transport holonomy; the ansatz's $3\times2\pi$ total coordinate
winding is separately accounted for.
\end{remark}
\status{published}
\source{P16:rem:holonomy_distinct}

\begin{construction}[Holonomy-compensated, arc-length-uniform frame]
\label{P16:constr:compensated_frame}
Add a compensating rotation about the tangent, distributed proportionally
to arc length:
\begin{equation}
  \varphi_{\mathrm{comp}}(t) = -H\,\frac{s(t)}{L_3},
\end{equation}
where $s(t)=\int_0^t|\Gamma'(u)|\,du$ is arc length and $L_3$ is the
total arc length. The corrected frame closes exactly.
\end{construction}
\status{published}
\source{P16:constr:compensated_frame}

\begin{proposition}[Frame closure]
\label{P16:prop:frame_closure}
The frame of Construction~\ref{P16:constr:compensated_frame} closes to within
$3\times10^{-6}$ degrees, consistent with exact closure up to
floating-point accumulation.
\end{proposition}
\status{published}
\source{P16:prop:frame_closure}

\begin{construction}[Per-strand poloidal phase]
\label{P16:constr:perstrand}
For a point $\mathbf{x}$ near strand $i$ (curve parameter $t_i$,
nearest point $\Gamma(t_i)$), define the poloidal angle $\chi_i(\mathbf{x})$
as the angle of $\mathbf{x}-\Gamma(t_i)$ in the $(N_1(t_i),N_2(t_i))$
plane of the compensated frame (Construction~\ref{P16:constr:compensated_frame}),
and set
\begin{equation}
  \Phi_i(\mathbf{x}) = \chi_i(\mathbf{x}) + 3t_i,
  \qquad
  Z_i(\mathbf{x}) = \bigl(\cos\tfrac{f_i}{2},\,\sin\tfrac{f_i}{2}\,e^{i\Phi_i}\bigr).
\end{equation}
At a crossing, combine $Z_1,Z_2$ by inverse-square-weighted complex
superposition as in Construction~\ref{P16:constr:s3lift}.
\end{construction}
\status{published}
\source{P16:constr:perstrand}

\begin{proposition}[$\chi$ is a genuine poloidal angle]
\label{P16:prop:chi_winds}
$\chi_i(\mathbf{x})$ winds exactly once around any small transverse loop
encircling the curve at fixed $t_i$; winding number confirmed as $1.0000$
by direct evaluation (script \texttt{perstrand\_winding\_test.py}).
\end{proposition}
\status{published}
\source{P16:prop:chi_winds}

\begin{proposition}[Construction~\ref{P16:constr:perstrand} is smooth at all three crossings]
\label{P16:prop:perstrand_smooth}
The field of Construction~\ref{P16:constr:perstrand} is continuous and
differentiable across the midline of all three crossing regions.
A resolution scan of local gradient energy
($N_{\mathrm{grid}}\in\{41,61,81,121\}$) gives
$g^2_{\max}\in\{56.2,58.4,59.3,60.0\}$, converging cleanly to a finite
value comparable to the chi-free construction's $\approx53$.
\end{proposition}
\status{published}
\source{P16:prop:perstrand_smooth}

\begin{theorem}[The corrected per-strand construction restores $\QH=3$]
\label{P16:thm:perstrand_charge_three}
The Whitehead-integral Hopf charge of the global per-strand construction
converges to $\QH=3$, ruling out $\QH=0,1,2$: the sequence
$1.052\to1.176\to1.796\to2.499$ at $N=64,70,120,240$ is monotonically
increasing for $N\geq64$ and already exceeds $2$ at $N=240$. The deficit
$3-\QH(N)$ fits a power law $\propto N^{-1.02}$, consistent with
$O(1/N)$ accuracy of a central-difference curl integral.
\end{theorem}
\status{published}
\source{P16:thm:perstrand_charge_three}

\begin{remark}[Status of the per-strand construction]
\label{P16:rem:perstrand_summary}
Construction~\ref{P16:constr:perstrand} combined with
Construction~\ref{P16:constr:compensated_frame} is the first ansatz in this
series to satisfy all three requirements simultaneously: smoothness at
every crossing (Proposition~\ref{P16:prop:perstrand_smooth}), bounded gradient
energy under refinement, and correct topological charge $\QH=3$
(Theorem~\ref{P16:thm:perstrand_charge_three}). This resolves the gap left
by Remark~\ref{P16:rem:s3lift_not_sufficient}: the per-strand poloidal winding
was the missing structural ingredient.
\end{remark}
\status{published}
\source{P16:rem:perstrand_summary}

%════════════════════════════════════════════════════════════════════
% Solver Functional Test
%════════════════════════════════════════════════════════════════════

\begin{proposition}[Global functional comparison at the solver's working resolution]
\label{P16:prop:global_functional_comparison}
At the solver's default grid ($N=64$, $h=0.26$, $C^*=2.4987$), the
three constructions compare as:
\begin{center}
\begin{tabular}{lrrrr}
\toprule
Construction & $K$ & $J_4$ & $E_{\mathrm{geom}}=K\cdot J_4$ \\
\midrule
A (original, discontinuous) & 2911.2 & 126.6 & 368,438 \\
B (chi-free $S^3$-lift, $\QH=0$) & 1510.4 & 76.8 & 115,960 \\
C (per-strand, $\QH=3$) & 1736.4 & 2181.3 & 3,787,624 \\
\bottomrule
\end{tabular}
\end{center}
The large $E_{\mathrm{geom}}$ of Construction~C at coarse resolution is
a resolution artifact, not a divergence.
\end{proposition}
\status{published}
\source{P16:prop:global_functional_comparison}

\begin{proposition}[The large $J_4$ is a resolution artifact, not a divergence]
\label{P16:prop:j4_resolution_artifact}
The $J_4$ excess of Construction~\ref{P16:constr:perstrand} ($98\%$ concentrated
within $\rho<1.0$) converges cleanly to a finite plateau under local
refinement near a crossing
($J_4\to453\to632\to729\to749\to729$ at $N=41,61,81,121,161$), in sharp
contrast to Construction~A, whose $J_4$ grows without bound under the same
refinement. The global excess matches a tiling estimate from the generic
tube-winding contribution ($2128.8$ predicted vs.\ $2181.3$ measured, $2.4\%$).
\end{proposition}
\status{published}
\source{P16:prop:j4_resolution_artifact}

\begin{remark}[Interpretation: generic resolution requirement]
\label{P16:rem:j4_resolution_interpretation}
The poloidal angle $\chi_i$ satisfies $|\nabla\chi_i|\sim1/\rho$ near
the curve, present at every point along the tube. The solver's default
grid ($h=0.26$) gives fewer than two points across the tube diameter,
under-resolving this winding everywhere. This is the standard resolution
requirement for any field with genuine angular winding around a thin core
(vortex lines, cosmic strings, monopole cores). Construction~A's divergence
under refinement is a genuine discontinuity, not a resolvable coarse-grid
effect; Construction~\ref{P16:constr:perstrand}'s excess is finite, resolvable,
and fully accounted for by the generic tube-winding mechanism.
\end{remark}
\status{published}
\source{P16:rem:j4_resolution_interpretation}

\begin{remark}[Practical recommendation for solver use]
\label{P16:rem:solver_practical_recommendation}
Using Construction~\ref{P16:constr:perstrand} as an actual solver initial
condition is not yet demonstrated: adequately resolving the tube core
requires $h\lesssim0.03$, which is not feasible on a uniform global grid
at the box sizes used in this series ($N^3\sim5\times10^8$ points
required, versus $N^3=2.6\times10^5$ at the solver's current default).
A non-uniform or adaptively refined mesh concentrated within one to two
tube radii of $T_{2,3}$ is the natural next step.
\end{remark}
\status{open}
\source{P16:rem:solver_practical_recommendation}

%════════════════════════════════════════════════════════════════════
% Gradient-Flow Landscape
%════════════════════════════════════════════════════════════════════

\begin{proposition}[Phase 1 energy removal]
\label{P16:prop:gf_phase1}
Starting from Construction~\ref{P16:constr:perstrand} at $N=144$, $h=0.075$,
gradient flow with $\mathrm{lr}=3\times10^{-4}$ reduces $E_{\mathrm{geom}}$
by a factor of $\approx2$ in 710 steps, with $\bar r$ remaining within
$0.018$ of $R_0=3$ throughout. The $K$-minimum at step $\approx711$
($K_{\min}=2313.9$; $K_{\min}\approx2294$ at $N=192$, $h=0.05$) marks
the point at which the initial excess gradient energy has been removed
and the field enters the nontrivial $K$--$J_4$ competition region.
\end{proposition}
\status{published}
\source{P16:prop:gf_phase1}
\residual{$K_{\min}\approx2294$ at $N=192$ (resolution-convergent)}

\begin{proposition}[Topology loss via cumulative drift]
\label{P16:prop:gf_topology_loss}
On the $N=144$, $h=0.075$ grid, gradient flow from
Construction~\ref{P16:constr:perstrand} does not converge to a $\QH=3$
energy minimum. The field instead drifts continuously toward the vacuum
($J_4\to0$, $\QH\to0$, $E_{\mathrm{geom}}\to0$) through slow,
cumulative small-angle rotations, each individually within the
discrete-topology-safe regime. The per-point rotation clamp prevents
catastrophic single-step topology changes but cannot prevent this
cumulative mechanism.
\end{proposition}
\status{published}
\source{P16:prop:gf_topology_loss}

\begin{remark}[Two-phase landscape and $(2+1)$ lobe splitting]
\label{P16:rem:gf_twophase}
Phase~1 (steps 1--710) operates within the $\QH=3$ sector, shedding
excess gradient energy. Phase~2 shows a $(2+1)$ lobe-splitting pattern
in the $N=192$ run: one outer lobe together with its adjacent midpoint
are depleted while two lobes retain their flux. This is the lightest
single string-breaking mode of the baryon: one colour flux tube is
released while the remaining two-colour diquark system stays bound.
The resolution-convergent $K$-minimum estimate is $K_{\min}\approx2294$
($N=192$, $h=0.05$), shifting by $1.0\%$ from the $N=144$ result
($K_{\min}=2313.9$), consistent with $O(h^2)$ grid error.
\end{remark}
\status{published}
\source{P16:rem:gf_twophase}

%════════════════════════════════════════════════════════════════════
% Arc-Segment Energy Partition
%════════════════════════════════════════════════════════════════════

\begin{definition}[Arc-segment boundaries]
\label{P16:def:arc_boundaries}
Define six distinguished parameter values on $T_{2,3}$ in monotone order:
\begin{align}
  t_{\mathrm{C}_n} &= \tfrac{\pi}{6}+\tfrac{2\pi n}{3}, \qquad
  t_{\mathrm{M}_n} = \tfrac{\pi}{3}+\tfrac{2\pi n}{3}, \quad n=0,1,2.
\end{align}
\textbf{B-segments} (crossing to midpoint, $[\mathrm{C}_n,\mathrm{M}_n]$):
arc-parameter width $\pi/6$ (30°).
\textbf{A-segments} (midpoint to crossing, $[\mathrm{M}_n,\mathrm{C}_{n+1}]$):
arc-parameter width $\pi/2$ (90°).
\end{definition}
\status{published}
\source{P16:def:arc_boundaries}

\begin{proposition}[$\mathbb{Z}_3$ symmetry of arc-segment energies]
\label{P16:prop:z3_symmetry}
For Construction~\ref{P16:constr:perstrand} on a cubic grid, the three
B-segment $J_4$ integrals agree to within $2.5\%$ ($N=64$) and $1.4\%$
($N=96$), and the three A-segment integrals agree to within $0.6\%$ and
$0.3\%$ respectively. The residual asymmetry decreases with resolution,
consistent with a grid discretisation artifact. The field partitions its
$J_4$ charge in a $\mathbb{Z}_3$-symmetric pattern to numerical precision.
\end{proposition}
\status{published}
\source{P16:prop:z3_symmetry}

\begin{proposition}[Geometric data of the three distinguished point types]
\label{P16:prop:geom_data}
For the trefoil $T_{2,3}$ with $R_0=3$, $r_0=0.874$:
\begin{enumerate}
\item \textbf{Crossings} ($t=t_{\mathrm{C}_n}$): curvature $\kappa=0.399$
  (global maximum), torsion $\tau\approx-0.17$, $S_{\mathrm{eff}}$ at
  maximum, $\rho_{J_4}$ at local minimum. All at $z=+r_0=+0.874$, $R_{xy}=R_0=3$.
\item \textbf{Midpoints} ($t=t_{\mathrm{M}_n}$): curvature $\kappa=0.026$
  (global minimum, factor~15 below crossings), torsion $|\tau|\approx11.3$
  (factor~67 above crossings), $S_{\mathrm{eff}}$ at minimum, colour
  weights exactly equal. All at $z=0$, $R_{xy}=R_0-r_0=2.126$.
\item \textbf{Outer-lobe distal points} ($t_{\mathrm{D}_n}=2\pi n/3$):
  curvature $\kappa=0.349$ (intermediate), torsion $\tau\approx0$ (local
  $\mathbb{Z}_2$ symmetry), $S_{\mathrm{eff}}$ intermediate. All at $z=0$,
  $R_{xy}=R_0+r_0=3.874$, each equidistant ($d=3.626$) from its two
  adjacent crossing vertices.
\item \textbf{Arc interiors}: A-segment bodies carry $\approx79\%$ of total
  $J_4$ at density $\rho_A\approx864$ per unit arc-parameter; B-segment
  bodies carry $\approx21\%$ at $\rho_B\approx700$. Energy density is
  genuinely higher on A-segments (ratio $\rho_B/\rho_A\approx0.81$, stable
  across resolutions).
\end{enumerate}
\end{proposition}
\status{published}
\source{P16:prop:geom_data}
\residual{exact (from $T_{2,3}$ geometry)}

\begin{remark}[Torsion contrast and field blindness at midpoints]
\label{P16:rem:torsion_contrast}
The torsion ratio $|\tau_{\mathrm{M}}|/|\tau_{\mathrm{C}}|\approx67$ is
dramatic: the midpoints are the sites of the most violent Frenet-frame
rotation on the entire curve, yet Paper~XV proves torsion makes zero
contribution to $S_{\mathrm{eff}}$ at all orders (Theorem~6.1 of Paper~XV).
The field is maximally blind to geometry exactly where the geometry is most
active. The midpoints are structurally weak loci: under-compressed
($\kappa_{\min}$), maximally diffuse colour, zero $S_{\mathrm{eff}}$
restoring force from torsion.
\end{remark}
\status{published}
\source{P16:rem:torsion_contrast}

\begin{remark}[$\rho_{J_4}$ anti-correlates with $S_{\mathrm{eff}}$ along the tube]
\label{P16:rem:rhoJ4_anticorrelation}
The arc-segment integrals establish an anti-correlation: $\rho_{J_4}$ is
lowest precisely where $S_{\mathrm{eff}}\propto\kappa^2$ is highest (at
the crossings), and the bulk of $J_4$ flux is concentrated in the A-segment
bodies where $\kappa$ is intermediate. The Hopf flux is repelled from
regions of high tube curvature; crossing fraying is a deficit signature,
not a concentration.
\end{remark}
\status{published}
\source{P16:rem:rhoJ4_anticorrelation}

\begin{remark}[Physical interpretation: 12-site fraying and decay geometry]
\label{P16:rem:arcseg_physics}
The $\rho_{J_4}$ field shows 12 distinct fraying sites in three categories
(6 crossing-associated, 3 midpoint torsion-twist, 3 outer-lobe distal),
each category internally related by $\mathbb{Z}_3$. If the three crossings
($z=+r_0$, $\kappa_{\max}$, pure colour, $\rho_{J_4}$ minimum) are
identified as quark colour vertices and the A-segment bodies as flux-tube
arms, the flux-tube energy is concentrated in the arm interiors and depleted
at the vertices, consistent with the standard QCD picture. The crossing
vertex acts as a topological lens: incoming A-segment flux is redirected
into the adjacent B-segment departure zone, the primary jet-ejection site.
\end{remark}
\status{published}
\source{P16:rem:arcseg_physics}

%════════════════════════════════════════════════════════════════════
% Trefoil Geometry: Frenet-Serret and Torsion — Paper XV
%════════════════════════════════════════════════════════════════════

\begin{theorem}[Exact torsion cancellation for symmetric profiles]
\label{P15:thm:torsion_cancels}
In the Frenet--Serret frame, the torsion $\tau(s)$ acts as a rigid
rotation $\chi\to\chi-\int_0^s\tau(s')\,ds'$ of the tube cross-section
coordinate. For any profile $f=f_0(\rho)$ depending only on $\rho$
(i.e.\ $\partial f/\partial\chi=0$), this rotation is invisible at
\emph{every order}: the torsion makes zero contribution to
$J_4^{(3)}/J_{2a}^{(3)}$ at all orders in $\rho\tau$.
\end{theorem}
\status{published}
\source{P15:thm:torsion_cancels}

\begin{remark}[Scope of the cancellation: leading order only]
\label{P15:rem:torsion_scope}
The cancellation is exact for the leading-order, $\chi$-independent
profile $f_0(\rho)$ alone. Once the curvature-sourced correction
$f_1(\rho)\cos\chi$ is included, the leading surviving contribution
is
\begin{equation}
  \bigl\langle\delta|\nabla f|^2\bigr\rangle_\chi
  = \tfrac{1}{2}\,f_1(\rho)^2\,\kappa(s)^2\,\tau(s)^2\,\rho^2
  + O(\rho^3),
  \label{P15:eq:torsion_correction}
\end{equation}
proportional to $\kappa^2\tau^2$, not $\tau^2$ alone.
Along the trefoil, $\langle\kappa^2\tau^2\rangle\approx0.0053$,
about $4\%$ of $\langle\kappa^2\rangle=0.120$.
This correction is developed in Paper~XVI.
\end{remark}
\status{published}
\source{P15:rem:torsion_scope}

\begin{corollary}[Bogomolny parameter from $C_3^*$]
\label{P15:cor:lambda3}
Combining Theorem~\ref{P15:thm:virial} ($\Kfb=2\vp\,\Ja$)
with the leading-order trefoil formula
$J_4^{(3)}/J_{2a}^{(3)}=3/(4C_3^{*2})+\langle\kappa^2\rangle/2$,
the leading-order Bogomolny parameter of the $\QH=3$ sector is
\begin{equation}
  \boxed{\lam_3
  = \frac{\Kfb}{J_4^{(3)}}
  = \frac{2\vp}{\dfrac{3}{4C_3^{*2}}+\dfrac{\langle\kappa^2\rangle}{2}},}
  \label{P15:eq:lambda3}
\end{equation}
at leading order in $\rho/R_0$.
If $\lam_3=\vp^6$, then $C_3^*=2.497$ at $O(\langle\kappa^2\rangle)$,
corrected to $C_3^*=2.5062$ at $O(\langle\kappa^4\rangle)$
(Proposition~\ref{P15:prop:kappa4}).
\end{corollary}
\status{published}
\source{P15:cor:lambda3}
\residual{exact}

%════════════════════════════════════════════════════════════════════
% Crossing Regions and Confinement — Paper XV
%════════════════════════════════════════════════════════════════════

\begin{proposition}[Crossing regions as inner wall]
\label{P15:prop:crossings}
The three crossing regions of the trefoil $T_{2,3}$ are the
$\QH=3$ analogue of the torus inner wall:
\begin{enumerate}[noitemsep,label=(\roman*)]
\item They are the loci of maximum curvature ($\kappa_{\max}=0.399$)
  and hence maximum Shafranov shift along the trefoil.
\item They carry enhanced condensate density by the same
  curvature-driven mechanism as the $3.6\times$ inner/outer ratio
  of $\QH=2$.
\item They are related by the $\mathbb{Z}_3$ rotational symmetry
  of the trefoil.
\end{enumerate}
\end{proposition}
\status{published}
\source{P15:prop:crossings}

\begin{remark}[Three concentrations: exact count]
\label{P15:rem:3or6}
The map $t\mapsto t+\pi$ combined with $(x,y,z)\mapsto(x,y,-z)$
is an exact isometry of $T_{2,3}$ that exchanges the over- and under-strands
at every crossing, forcing $\kappa_{\mathrm{over}}=\kappa_{\mathrm{under}}=0.3991$
for all three crossings.
\begin{equation}
  \kappa(t_k)=\kappa(t_k+\pi)=0.3991.
  \label{P15:eq:kappa_equal}
\end{equation}
Each crossing region is therefore a single energy concentration
with $\mathbb{Z}_2$-symmetric over/under contributions.
The count is exactly $3$: one per crossing.
\end{remark}
\status{published}
\source{P15:rem:3or6}

\begin{proposition}[Confinement from boundary conditions]
\label{P15:prop:confinement}
Colour confinement is encoded in the profile boundary conditions
\begin{equation}
  f(0) = \pi,\qquad f(\infty) = 0.
  \label{P15:eq:BC}
\end{equation}
An isolated open flux strand requires $f$ to interpolate $0\to\pi\to0$
along its length; the gradient energy of such a configuration diverges
linearly with strand length.
Only closed configurations have finite energy, requiring the three
trefoil tubes to link into the closed knot $T_{2,3}$.
\end{proposition}
\status{published}
\source{P15:prop:confinement}

\begin{remark}[Absence of a unique Bogomolny parameter and confinement]
\label{P15:rem:no_lambda3}
Extensive numerical gradient flow on $80^3$ grids consistently shows
$\Kfb/J_4$ traversing $\vp^6$ without converging there, due to the
three-lobed energy landscape from the crossing-region concentrations.
The bound $\lam_3>\vp^5=11.09$ is established numerically.
The appropriate characterisation is through the global invariants
$N_3=3$ (Theorem~\ref{P15:thm:Qgroup3}) and $\Kfb=2\vp\Ja$
(Theorem~\ref{P15:thm:virial}), rather than a local saddle value.
Since $\lam_3=\vp^6$ is proved (Theorem~\ref{P15:thm:sector_assignment}),
this gives $C_3^*=2.5062$ analytically (Proposition~\ref{P15:prop:kappa4}).
\end{remark}
\status{published}
\source{P15:rem:no_lambda3}

\begin{proposition}[Y-junction geometry: the Fermat--Steiner energy balance]
\label{P15:prop:yjunction}
The three crossing midpoints
$\mathbf{m}_k = \tfrac{1}{2}[\Gamma(t_k)+\Gamma(t_k+\pi)]$,
$k=1,2,3$, satisfy
\begin{equation}
  |\mathbf{m}_k| = R_0,\qquad
  |\mathbf{m}_k - \mathbf{m}_j| = R_0\sqrt{3}\quad(j\neq k),\qquad
  \sum_{k=1}^3 \mathbf{m}_k = \mathbf{0}.
  \label{P15:eq:yjunction}
\end{equation}
The three midpoints form a regular equilateral triangle in the $z=0$
plane; the origin is the Fermat--Steiner point with three equal arms
of length $R_0$.
\end{proposition}
\status{published}
\source{P15:prop:yjunction}
\residual{exact}

\begin{remark}[The Y-junction as the missing energy term]
\label{P15:rem:yjunction_energy}
The full baryon energy is
\begin{equation}
  E_{\mathrm{baryon}} = \underbrace{K_{\mathrm{fb}}\times J_4}_{\text{profile}}
  + \underbrace{\sigma\,\cdot\,3R_0}_{\text{string tension}}
  + \underbrace{E_\cup}_{\text{junction}},
  \label{P15:eq:baryon_energy}
\end{equation}
with total Y-junction string length $3R_0 = N_3^2 = 9$.
Numerical flows minimised only the profile term, missing the
string-tension contribution. The profile saddle (determining $C_3^*$
and $\lam_3$) is unchanged since the missing term is $C^*$-independent.
\end{remark}
\status{published}
\source{P15:rem:yjunction_energy}

%════════════════════════════════════════════════════════════════════
% Normalisation Identity and Equivalence — Paper XV
%════════════════════════════════════════════════════════════════════

\begin{theorem}[Normalisation identity]
\label{P15:thm:norm3}
Given Theorem~\ref{P15:thm:virial} ($\Kfb=2\vp\Ja$), the sector energy
formula $\Efb^{(3)}=2N_3^2/\vp^{17}$, the thermodynamic axiom
$\Efb^{(3)}/V_3=\rCMB=10$, and $\lam_3=\vp^6$:
\begin{equation}
  \boxed{\vp^9\sqrt{\Ja^{(3)}\Jfour^{(3)}}
  = 3 = Q_{\mathrm{group}}^{(3)}.}
  \label{P15:eq:norm3}
\end{equation}
\end{theorem}
\status{published}
\source{P15:thm:norm3}
\residual{exact}

\begin{corollary}[Condensate volume as derived quantity]
\label{P15:cor:V3}
Under the same hypothesis $\lam_3=\vp^6$,
the trefoil condensate volume is
\begin{equation}
  V_3 = \frac{N_3^2}{5\vp^{17}} = \frac{9}{5\vp^{17}}.
  \label{P15:eq:V3}
\end{equation}
This is a derived consequence of $\lam_3=\vp^6$, not an independent
assumption.
\end{corollary}
\status{published}
\source{P15:cor:V3}
\residual{exact}

\begin{theorem}[Logical equivalence of the two closure conditions]
\label{P15:thm:equivalence}
The following two statements are logically equivalent, given the
thermodynamic axiom and energy formula:
\begin{enumerate}[noitemsep,label=(\Roman*)]
\item $\lam_3 = \vp^6$ (sector-independent Bogomolny parameter).
\item $\vp^9\sqrt{\Ja^{(3)}\Jfour^{(3)}} = 3$ (normalisation identity).
\end{enumerate}
Either implies $V_3=9/(5\vp^{17})$ as a corollary.
\end{theorem}
\status{published}
\source{P15:thm:equivalence}
\residual{exact}

\begin{proposition}[$O(\langle\kappa^4\rangle)$ correction via Beta functions]
\label{P15:prop:kappa4}
The profile integrals reduce to Beta functions via $t=\hat\rho^2$:
$I_n = 2^{2n+1}B(n+1,n+1)$, giving $I_2=16/15$.
The $O(\rho^4\kappa^4)$ curvature correction to $r_3$ is
\begin{equation}
  \delta r_3^{(\kappa^4)}
  = \frac{9\,\langle\kappa^4\rangle}{4\,C_3^{*4}}
  = 0.000938 \quad (0.520\%\text{ of }r_3),
  \label{P15:eq:dr_kappa4_num}
\end{equation}
using $\langle\kappa^4\rangle=0.0162$ for $T_{2,3}$.
This shifts $C_3^*$ from $2.497$ to
\begin{equation}
  \boxed{C_3^* = 2.5062 \quad\bigl(O(\langle\kappa^4\rangle)\text{ corrected}\bigr).}
  \label{P15:eq:C3star_kappa4}
\end{equation}
\end{proposition}
\status{published}
\source{P15:prop:kappa4}
\residual{exact}

%════════════════════════════════════════════════════════════════════
% Condensate Universality and Simple-Current Sector — Paper XV
%════════════════════════════════════════════════════════════════════

\begin{proposition}[Condensate universality of $\lam_3$]
\label{P15:prop:lambda3}
If the $\QH=3$ density-feedback functional admits a topology-preserving
saddle with virial fixed point $V^*=\vp$, and if the EL saddle profile
satisfies $r_3 = J_4^{(3)}/J_{2a}^{(3)} = 2/\vp^5$, then
$\lam_3=\vp^6$.
\end{proposition}
\status{published}
\source{P15:prop:lambda3}
\residual{exact}

\begin{proposition}[Exact structure of $r_3$ corrections]
\label{P15:prop:r3_exact}
For any axially-symmetric profile $f_0(\rho)$:
\begin{enumerate}[label=\textup{(\roman*)}]
\item $J_{2a}^{(3)}$ receives no curvature correction from the
  Frenet--Serret Jacobian: $\delta J_{2a}^\kappa = 0$ exactly.
\item The leading curvature correction to $J_4^{(3)}$ is
  $\langle F_{s\chi}^2\rangle_\chi = \tfrac12\kappa^2\sin^4\!f_0(f_0')^2$.
\item The ratio has the exact expansion
  \begin{equation}
    r_3 = \frac{3}{4C_3^{*2}}
    + \frac{\langle\kappa^2\rangle_s}{2}\cdot\Gamma(C_3^*)
    + O(\langle\kappa^4\rangle_s).
    \label{P15:eq:r3_exact_expansion}
  \end{equation}
\end{enumerate}
\end{proposition}
\status{published}
\source{P15:prop:r3_exact}

\begin{lemma}[Fusion-count ratio of sector targets]
\label{P15:lem:fusion_ratio}
The target ratios for the two sectors satisfy
\begin{equation}
  \frac{r_3^*}{r_2^*}
  = \frac{2/\vp^5}{2^{4/3}/\vp^5}
  = \frac{1}{2^{1/3}}
  = N_{\mathrm{fus}}^{-1/k},
  \label{P15:eq:fusion_ratio}
\end{equation}
where $N_{\mathrm{fus}}=2$ (fusion count) and $k=3$ (WZW level).
\end{lemma}
\status{published}
\source{P15:lem:fusion_ratio}
\residual{exact}

\begin{theorem}[Simple-current sector assignment]
% Note: label retains historical 'conj:' prefix; type promoted to theorem.
\label{P15:conj:sector}
The $\QH=3$ trefoil Hopfion belongs to the $j=k/2=3/2$ simple-current
sector of the $\mathrm{SU}(2)_3$ condensate WZW model.
Three independent proofs are given (Theorem~\ref{P15:thm:sector_assignment},
Proposition~\ref{P15:prop:verlinde_unique}, and the $\mathbb{Z}_2$ parity
filter of Section~\ref{sec:z2_parity}).
\end{theorem}
\status{published}
\source{P15:conj:sector (promoted from conjecture to theorem)}

\begin{theorem}[$\lam_3=\vp^6$]
\label{P15:thm:lambda3_sc}
The $\QH=3$ trefoil is in the $j=3/2$ sector of $\mathrm{SU}(2)_3$
(Theorem~\ref{P15:conj:sector}), so $\lam_3 = \vp^6$:
\begin{equation}
  \boxed{\lam_3 = \frac{\lam}{s_{\mathrm{opt}}^{2k}} = \frac{\vp^6}{1} = \vp^6.}
  \label{P15:eq:lambda3_proved}
\end{equation}
\end{theorem}
\status{published}
\source{P15:thm:lambda3_sc}
\residual{exact}

\begin{corollary}[All results now unconditional]
\label{P15:cor:sc_consequences}
\begin{enumerate}[noitemsep,label=\textup{(\roman*)}]
\item $r_3 = 2/\vp^5$.
\item $\vp^9\sqrt{\Ja^{(3)}\Jfour^{(3)}}=3$.
\item $V_3 = 9/(5\vp^{17})$.
\item $C_3^* = 2.5062$ (at $O(\langle\kappa^4\rangle)$).
\end{enumerate}
All four follow from $\lam_3=\vp^6$ via the algebraic chain of
Theorem~\ref{P15:thm:virial}--Corollary~\ref{P15:cor:V3}.
\end{corollary}
\status{published}
\source{P15:cor:sc_consequences}
\residual{exact}

%════════════════════════════════════════════════════════════════════
% WRT Invariants — Paper XV
%════════════════════════════════════════════════════════════════════

\begin{proposition}[Writhe-3 framing phase of the simple current]
\label{P15:prop:wrt_phase}
For $j=3/2$ at $k=3$, the writhe-3 framing phase is
\begin{equation}
  \theta_{3/2}^{\,3}
  = e^{2\pi i\cdot 3h_{3/2}}
  = e^{2\pi i\cdot 9/4}
  = e^{2\pi i/4}
  = i.
  \label{P15:eq:theta32_cubed}
\end{equation}
This is a primitive $4$th root of unity, the smallest non-trivial order
among all sectors at writhe $w=3$.
\end{proposition}
\status{published}
\source{P15:prop:wrt_phase}
\residual{exact}

\begin{proposition}[String tension minimisation by the simple current]
\label{P15:prop:wrt_dim}
For a closed knot, the WZW string tension per unit length is proportional
to $\dim_q(j)$. Among non-vacuum sectors ($j\neq0$) at $k=3$,
the minimum is achieved uniquely by $j=3/2$:
\begin{equation}
  \dim_q\!\bigl(\tfrac{3}{2};\,k{=}3\bigr) = 1
  \;<\; \vp \;=\; \dim_q\!\bigl(\tfrac{1}{2};\,k{=}3\bigr)
        = \dim_q\!\bigl(1;\,k{=}3\bigr).
  \label{P15:eq:dim_min}
\end{equation}
\end{proposition}
\status{published}
\source{P15:prop:wrt_dim}
\residual{exact}

\begin{proposition}[Pure-phase WRT contribution of $j=3/2$]
\label{P15:prop:wrt_pure}
The $j=3/2$ sector's contribution to $\tau_3(T_{2,3})$ is
\begin{equation}
  \frac{S_{0,\,3/2}}{S_{0,0}}\,\theta_{3/2}^3
  = 1\cdot i = i,
  \label{P15:eq:wrt_pure}
\end{equation}
which has unit modulus. Every other non-vacuum sector contributes
with modulus $\vp>1$.
\end{proposition}
\status{published}
\source{P15:prop:wrt_pure}
\residual{exact}

%════════════════════════════════════════════════════════════════════
% Proof of the Sector Assignment — Paper XV
%════════════════════════════════════════════════════════════════════

\begin{lemma}[Triple-fusion decomposition at $k=3$]
\label{P15:lem:triple_fusion}
In the $\mathrm{SU}(2)_3$ WZW model, the triple fusion of the
fundamental representation decomposes as
\begin{equation}
  j_{\frac12}\otimes j_{\frac12}\otimes j_{\frac12}
  = j_{\frac12}^{(\times 2)} \oplus j_{\frac32}^{(\times 1)}.
  \label{P15:eq:triple_fusion}
\end{equation}
The $j=1/2$ channel has multiplicity $2$ and the $j=3/2$ channel
has multiplicity $1$.
\end{lemma}
\status{published}
\source{P15:lem:triple_fusion}

\begin{theorem}[Sector assignment: $\QH=3\leftrightarrow j=3/2$]
\label{P15:thm:sector_assignment}
Under the Witten bosonization dictionary, the $\QH=3$ Hopfion occupies
the $j=3/2$ sector of the $\mathrm{SU}(2)_3$ condensate WZW model.
The proof proceeds in four steps: (1)~each $\QH=1$ quantum carries
$j=1/2$ (Paper~III Prop.~10.6); (2)~$\QH=3$ lies in
$j_{1/2}^{\otimes3}=j_{1/2}^{(\times2)}\oplus j_{3/2}^{(\times1)}$
(Lemma~\ref{P15:lem:triple_fusion}); (3)~$\pi_3(S^2)=\mathbb{Z}$ gives
topological non-degeneracy; (4)~the multiplicity-2 channel $j=1/2$
is excluded, leaving $j=3/2$ uniquely.
\end{theorem}
\status{published}
\source{P15:thm:sector_assignment}

\begin{corollary}[$\lam_3=\vp^6$: is closed]
\label{P15:cor:lambda3_proved}
The Bogomolny parameter of the $\QH=3$ sector is $\lam_3=\vp^6$.
By Theorem~\ref{P15:thm:equivalence}, all three equivalent conditions
follow simultaneously: $\lam_3=\vp^6$, $\vp^9\sqrt{\Ja^{(3)}\Jfour^{(3)}}=3$,
and $r_3=2/\vp^5$ at the EL saddle.
\end{corollary}
\status{published}
\source{P15:cor:lambda3_proved}
\residual{exact}

\begin{remark}[Logical status of Step 1]
\label{P15:rem:step1_status}
The identification ``each $\QH=1$ quantum carries $j=\frac12$'' is the
explicit content of Paper~III Proposition~10.6 and is not an additional
postulate of Paper~XV. Steps~2--4 of Theorem~\ref{P15:thm:sector_assignment}
are independent of this identification.
\end{remark}
\status{published}
\source{P15:rem:step1_status}

\begin{proposition}[Verlinde ratio uniquely selects $j=3/2$]
\label{P15:prop:verlinde_unique}
Among all primaries of $\mathrm{SU}(2)_3$,
$\bigl|S_{j,\,1/2}/S_{j,\,0}\bigr|=\vp$ if and only if
$j\in\{0,\,\frac32\}$:
\begin{equation}
  \frac{S_{j,\,1/2}}{S_{j,\,0}}
  = 2\cos\!\frac{(2j{+}1)\pi}{5},
  \label{P15:eq:ratio_formula}
\end{equation}
giving $+\vp,\,+1/\vp,\,-1/\vp,\,-\vp$ for $j=0,\tfrac12,1,\tfrac32$.
Since $j=0$ is $\QH=2$ and the sector-independent IVT requires
$|V^*|=\vp$, the $\QH=3$ sector must be $j=\frac32$.
\end{proposition}
\status{published}
\source{P15:prop:verlinde_unique}
\residual{exact}

\begin{remark}[Three independent arguments]
\label{P15:rem:three_arguments}
The sector assignment $\QH=3\leftrightarrow j=\frac32$ is supported
by three independent proofs:
\begin{enumerate}[noitemsep,label=\textup{(\roman*)}]
\item \emph{Multiplicity + topological uniqueness}
  (Theorem~\ref{P15:thm:sector_assignment}).
\item \emph{Verlinde uniqueness}
  (Proposition~\ref{P15:prop:verlinde_unique}).
\item \emph{$\mathbb{Z}_2$ parity filter}: $\QH=3$ (odd)
  maps to $\mathbb{Z}_2$-odd sectors $j\in\{\frac12,\frac32\}$;
  multiplicity then forces $j=\frac32$.
\end{enumerate}
None of the three arguments is circular with the others.
\end{remark}
\status{published}
\source{P15:rem:three_arguments}
